% Generated draft: requires mathematical and source review.
\documentclass[11pt]{article}
\usepackage[margin=25mm]{geometry}
\usepackage{fontspec}
\setmainfont{Latin Modern Roman}
\usepackage{amsmath,amssymb,mathtools,unicode-math}
\setmathfont{Latin Modern Math}
\usepackage{enumitem,needspace,xurl,hyperref}
\setlength{\parindent}{0pt}
\setlength{\parskip}{5pt}
\setlength{\emergencystretch}{3em}
\newcommand{\N}{\mathbb{N}}
\newcommand{\Z}{\mathbb{Z}}
\newcommand{\Q}{\mathbb{Q}}
\newcommand{\R}{\mathbb{R}}
\newcommand{\C}{\mathbb{C}}
\newcommand{\F}{\mathbb{F}}
\newcommand{\A}{\mathbb{A}}
\newcommand{\PP}{\mathbb{P}}
\newcommand{\E}{\mathbb{E}}
\newcommand{\eps}{\varepsilon}
\newcommand{\dd}{\,\mathrm{d}}
\begin{document}
\section*{UnsolvedMath ID: 701530}
\textbf{Source problem code:} catalogue entry
\subsection{Definitions and mathematical statement}
Let \(G\) be a finite solvable group. A Carter subgroup of \(G\) is a nilpotent subgroup \(C\leq G\) satisfying \(N_G(C)=C\), where \(N_G(C)=\{g\in G:g^{-1}Cg=C\}\). Write
\[
|C|=\prod_{p\mid |C|}p^{a_p},\qquad \ell(C)=\sum_{p\mid |C|}a_p,
\]
with \(\ell(1)=0\). Define the nilpotent length \(h(G)\), also called the Fitting height, to be the least integer \(r\geq 0\) for which there exists a series
\[
1=G_0\leq G_1\leq\cdots\leq G_r=G
\]
whose terms are normal in \(G\) and whose factors \(G_i/G_{i-1}\) are nilpotent. In particular, \(h(1)=0\).

The problem asks whether there are constants \(a,b\geq 0\), independent of \(G\) and \(C\), such that
\[
\forall\,G\text{ finite solvable}\quad\forall\,C\leq G\text{ a Carter subgroup},\qquad h(G)\leq a\ell(C)+b.
\]
If a linear bound is unavailable, does there at least exist a polynomial \(P\in\R[x]\), nonnegative on the nonnegative integers and independent of \(G\) and \(C\), such that
\[
\forall\,G\text{ finite solvable}\quad\forall\,C\leq G\text{ a Carter subgroup},\qquad h(G)\leq P(\ell(C))?
\]
Here ``linear'' is interpreted as an affine upper bound; the supplied record does not specify a convention for the constant term.

The supplied statement attributes to E. C. Dade an exponential upper bound in \(\ell(C)\). The supplied editors' comment attributes to E. Jabara a quadratic upper bound in the special case where, under the preceding hypotheses, \(G=H\rtimes C\), \(C\) is cyclic, and
\[
C_H(C)=\{h\in H:hc=ch\text{ for every }c\in C\}=1.
\]
Here \(G=H\rtimes C\) means that \(H\trianglelefteq G\), \(G=HC\), and \(H\cap C=1\). These literature assertions are supplied source qualifications and have not been independently verified.
\subsection{Short English statement}
Is the nilpotent length of every finite solvable group bounded by a universal linear function, or at least a universal polynomial, of the number of prime factors of the order of a Carter subgroup, counted with multiplicity? The supplied record reports an exponential bound in general and a quadratic bound for a specified cyclic complement case.
\subsection{Sources}
\begin{itemize}
\item Local UnsolvedMath registry, permanent ID \(701530\), exact source code \texttt{catalogue entry}, title \emph{Kourovka Notebook Problem 18.102}. The supplied registry labels the problem as open; this label is not an independently verified assessment.
\item \emph{The Kourovka Notebook (Unsolved Problems in Group Theory)}, original problem 18.102, 18th issue (2014); supplied background also identifies the 21st issue edition (2026), \texttt{arXiv:1401.0300v48}. The registry describes the collection as edited by E. I. Khukhro and V. D. Mazurov. Proposed by A. Turull in 2014; the statement separately attributes the problem to E. C. Dade.
\item E. C. Dade, \emph{Illinois J. Math.} \textbf{13} (1969), 449--514. Supplied citation for the exponential upper bound; no article title is supplied.
\item E. Jabara, \emph{J. Algebra} \textbf{487} (2017), 161--172. Supplied editors' citation for the quadratic bound when \(G=H\rtimes C\), \(C\) is cyclic, and \(C_H(C)=1\); no article title is supplied.
\end{itemize}
\end{document}
